\section{结论}
\label{sec:conclusion}

本文介绍了 WhatyTerm——一个面向命令行编码代理的无人值守终端编排系统——并以一次真实部署的实测数据为核心，报告了其设计意图与实际行为之间的落差。系统在设计上具备规则优先、LLM 兜底的分层决策，三通道感知，以及三层嵌套自主循环等机制；但跨三天、四会话、46 次决策的日志分析表明：全部决策 100\% 由规则层做出、LLM 兜底从未触发，87\% 的动作只是机械“继续”，自主引擎无一真正交付。一个本应智能编排自主开发的系统，退化成了“自动点继续”的宏。

我们将根因定位到规则层“检测到空闲提示符即无条件继续”的判定路径，并将其抽象为“廉价层垄断”反模式：当廉价规则层对高频状态总能给出自信的、语法驱动的答案，且缺乏动作有效性反馈时，昂贵的推理层会被系统性饿死，系统稳定收敛到“高效率、零有效性”的沉默空转。我们进一步指出，常规效率指标会掩盖这一退化，唯有引入任务推进度等有效性指标才能将其暴露。

本文的立场是：\emph{一个诚实的负面结果，连同其可复现的证据与代码级根因，比一份夸耀设计优越的报告更有价值}。我们希望“廉价层垄断”这一反模式、以及“效率指标须与有效性指标并列”的评估主张，能为自主代理 harness 的设计者提供一面镜子：真正的挑战不在于让系统跑起来，而在于让它在跑的时候确实在做有意义的事。未来工作包括实现并验证本文提出的移交条件与有效性反馈机制，并在更大规模的部署上检验退化是否被消除。
